\newpage
\section{拓扑与几何}
\subsection{拓扑学基本知识}
\begin{tcolorbox}[colback=KuangNei,colframe=BianKuang,title=\textbf{拓扑空间},sharpish corners]
    设$X$是一个集合,$\mathcal{T} \subset \mathcal{P} (X)$,如果$\ttt$满足下面三条性质:\\
    1. $\emptyset \in \mathcal{T},X\in \mathcal{T}$;\\
    2. $\forall A_i \in \mathcal{T} (i\in I)$,则$\bigcup_{i\in I}A_i \in \mathcal{T}$;\\
    3. $\forall A_i \in \mathcal{T} (i=1,2)$,则$A_1\cap A_2 \in \mathcal{T}$.
\end{tcolorbox}
就称$\ttt$为$X$上的一个拓扑,又称$\ttt$中的元素为$X$上的开集.当$X$装备了拓扑结构$\ttt$之后,称$(X,\ttt)$为拓扑空间.\par
\textbf{定义}$^\clubsuit $:设$(X,\ttt)$为一拓扑空间,$E\subset X$是闭集当且仅当$X-E$是开集;事实上,闭集还可以定义为满足$E=E'$的$E$为闭集.\par
设$(X,\ttt)$为一拓扑空间,$S$是$X$的一个子集,可以在$S$上定义一个与$\ttt$相容的拓扑
$$
\ttt _{S}=\{O\cap S;O\in\ttt\}
$$
称$\ttt _{S}$是$\ttt$在$S$上的诱导拓扑,$(S,\ttt _{S})$是$(X,\ttt)$的拓扑子空间.\par


\subsection{微分几何}
\subsubsection{拓扑基本知识}
\begin{Definition}{拓扑空间}
    设$X$是一个集合,$\mathcal{T} \subset \mathcal{P} (X)$,如果$\ttt$满足下面三条性质:\par
    1. $\emptyset \in \mathcal{T},X\in \mathcal{T}$;\par
    2. $\forall A_i \in \mathcal{T} (i\in I)$,则$\bigcup_{i\in I}A_i \in \mathcal{T}$;\par
    3. $\forall A_i \in \mathcal{T} (i=1,2)$,则$A_1\cap A_2 \in \mathcal{T}$.\par
    就称$\ttt$为$X$上的一个拓扑,又称$\ttt$中的元素为$X$上的开集.当$X$装备了拓扑结构$\ttt$之后,称$(X,\ttt)$为拓扑空间
\end{Definition}
我们还需要邻域的概念.
\begin{Definition}{邻域}
    设$X$是一个拓扑空间,$x\in A\subset X$.若存在$X$的开集$U$,使得$x\in U\subset A\subset X$,则称$A$是$x$的一个邻域.若$A$就是$X$的一个开集,则称$A$是$x$的一个开邻域.
\end{Definition}
\begin{Definition}{$T_2$分离性}
    设$X$是一个拓扑空间.若对于任意不同的两点$x,y\in X$,总存在$x$的邻域$U$和$y$的邻域$V$,使得$U\cap V=\emptyset$,则称$X$为$T_2$空间.
\end{Definition}
我们也称$T_2$空间为Housdorff空间.\par
一般来讲,我们无法完全确定一个空间的拓扑,因此,一个想法是从一部分开集来生成拓扑.这就引入了拓扑基的概念.\par
\begin{Definition}{拓扑基}
    设$X$是一个集合,$\mathcal{B}$是由$X$的若干非空子集构成的集合.若$\mathcal{B}$满足条件:\\
    1. 对于任意$x\in X$,存在$B\in \mathcal{B}$,使得$x\in B\in\mathcal{B}$;\\
    2. 如果$B_1,B_2\in\mathcal{B}$且$x\in B_1\cap B_2$,则存在$B\in\mathcal{B}$,使得$x\in B\subset B_1\cap B_2$.\\
    则称$\mathcal{B}$是$X$上的一个拓扑基.
\end{Definition}
如果令
$$
\mathcal{T}=\left\{U=\bigcup_{B\in \mathcal{U}}B:\mathcal{U}\subset\mathcal{B}\right\}
$$
则$\mathcal{T}$就是$X$上的拓扑.\par
对于有可数拓扑基的拓扑空间$X$,我们称之为$A_2$空间或者第二可数空间.
\subsubsection{微分流形}
\textbf{定义}$^\clubsuit $:设$M$是一个具有$A_2,T_2$性质的拓扑空间,如果存在$M$的开覆盖$\{U_{\alpha}\}_{\alpha\in\Gamma}$以及相应的连续映射族$\varphi_{\alpha}:U_{\alpha}\to\mathbb{R}^n$,使得\par
(i) $\varphi_{\alpha}:U_{\alpha}\to \varphi_{\alpha}(U_{\alpha})\subset\mathbb{R}^n$为从$U_{\alpha}$到欧氏空间$\varphi_{\alpha}(U_{\alpha})$上的同胚;\par
(ii) 当$U_{\alpha}\cap U_{\beta}\neq \emptyset$时,如下的转换映射
$$
\varphi_{\beta}\circ \varphi_{\alpha}^{-1}:\varphi_{\alpha}(U_{\alpha}\cap U_{\beta})\to \varphi_{\beta}(U_{\alpha}\cap U_{\beta})
$$
为$C^r(r\geqslant 1)$映射,则称$M$为$C^r$流形.\par
\begin{center}
\tikzset{
    manifold/.style={line width=0.8pt},
    chart/.style={line width=0.6pt, fill=gray!10, fill opacity=0.5},
    % 映射箭头：虚线 + 箭头
    map_arrow/.style={->, >=stealth, line width=0.5pt, dashed},
    % 装饰曲线：虚线，不带箭头
    deco_line/.style={line width=0.4pt, dashed, draw=gray!80}
}

\begin{tikzpicture}[x=0.75pt,y=0.75pt,yscale=-1,xscale=1]

% --- 1. 绘制流形 M ---
\draw[manifold] (181,89) .. controls (199,67) and (436,1.5) .. (446,35.5) 
    .. controls (456,69.5) and (467,142.5) .. (431,171.5) 
    .. controls (395,200.5) and (214,203.5) .. (194,173.5) 
    .. controls (174,143.5) and (163,111) .. (181,89) -- cycle ;
\node at (425, 50) {$M$};

% --- 2. 绘制坐标邻域 U_alpha 和 U_beta ---
\draw[chart] (244,98) .. controls (264,88) and (310,77.5) .. (338,101.5) 
    .. controls (366,125.5) and (323,159.5) .. (316,166.5) 
    .. controls (309,173.5) and (264,188) .. (244,158) 
    .. controls (224,128) and (224,108) .. (244,98) -- cycle ;

\draw[chart] (325,91) .. controls (351,80.5) and (407,81.5) .. (408,110.5) 
    .. controls (409,139.5) and (401,131.5) .. (401,154.5) 
    .. controls (401,177.5) and (367,168.5) .. (329,150.5) 
    .. controls (291,132.5) and (299,101.5) .. (325,91) -- cycle ;

\node at (275, 126) {$U_{\alpha}$};
\node at (375, 122) {$U_{\beta}$};

% --- 3. 绘制欧几里得空间中的像 (包含内部装饰虚线) ---

% 左侧像 phi_alpha(U_alpha)
\begin{scope}
    % 先绘制椭圆
    \draw[chart] (217,320.25) ellipse [x radius=59, y radius=50.25];
    % 裁剪：接下来的绘图只会在这个椭圆内显示
    \clip (217,320.25) ellipse [x radius=59, y radius=50.25];
    % 绘制内部装饰虚线
    \draw[deco_line] (238,367.5) .. controls (233,337.5) and (227,314.5) .. (253,279.5);
\end{scope}
\node at (217, 395) {$\varphi_{\alpha}(U_{\alpha})$};

% 右侧像 phi_beta(U_beta)
\begin{scope}
    % 先绘制椭圆
    \draw[chart] (453.5,318.25) ellipse [x radius=56.5, y radius=49.25];
    % 裁剪
    \clip (453.5,318.25) ellipse [x radius=56.5, y radius=49.25];
    % 绘制内部装饰虚线
    \draw[deco_line] (426,362.5) .. controls (466,332.5) and (431,296.5) .. (413,285.5);
\end{scope}
\node at (453.5, 397) {$\varphi_{\beta}(U_{\beta})$};

% --- 4. 绘制所有映射箭头 (全部使用虚线) ---

% 映射 phi
\draw[map_arrow] (263,179.5) -- (238,268.5) node[midway, left] {$\varphi_{\alpha}$};
\draw[map_arrow] (395,175.5) -- (438,265.5) node[midway, right] {$\varphi_{\beta}$};

% 逆映射 phi^-1
\draw[map_arrow] (262,271.5) -- (287,185.5) node[midway, right] {$\varphi_{\alpha}^{-1}$};
\draw[map_arrow] (412,272.5) -- (369,180.5) node[midway, left] {$\varphi_{\beta}^{-1}$};

% 转换映射
\draw[map_arrow] (283,320.5) -- (393,321.5) 
    node[midway, below=5pt] {$\varphi_{\beta} \circ \varphi_{\alpha}^{-1}$};

\end{tikzpicture}
\end{center}
我们称$\{(U_{\alpha},\varphi_{\alpha})\}$为$M$的局部坐标覆盖,$(U_{\alpha},\varphi_{\alpha})$为一个局部坐标系,$U_{\alpha}$为局部坐标邻域,$\varphi_{\alpha}$为局部坐标映射.设$p\in U_{\alpha}$,记$x^i(p)$为$\varphi_{\alpha}(p)$的第$i$个欧式坐标分量,则称$x^i:U_{\alpha}\to \mathbb{R}$为局部坐标系$(U_{\alpha},\varphi_{\alpha})$的第$i$个坐标函数,有时也称$\{x^i\}=\{x^1,x^2,\cdots,x^n\}$为$p$附近的局部坐标.\par
下面我们定义两个流形之间的映射.
\begin{Definition}{$C^{k}$映射}
    设$f:M\to N$为两个$C^r$微分流形之间的连续映射.如果任给$p\in M$和$q=f(p)\in N$附近的局部坐标系$(V,\psi )$,均存在$p$附近的局部坐标系$(U,\varphi)$,使得$f(U)\subset V$,且$f$在这两个局部坐标系下的局部表示
    $$
    \psi\circ f\circ \varphi^{-1}:\varphi(U)\to \psi(V)
    $$
    为$C^{k}(k\leqslant r)$映射,则称$f$为流形$M,N$之间的$C^{k}$映射.$C^{k}$映射的全体记为$C^{k}(M,N)$.
\end{Definition}
\begin{center}
    \tikzset{
    manifold/.style={line width=0.65pt},
    chart/.style={line width=0.65pt, fill=blue!5, fill opacity=0.4},
    % 映射箭头：虚线、stealth风格、缩短间距
    map_arrow/.style={->, >=stealth, dashed, line width=0.5pt, shorten >=3pt, shorten <=3pt},
    label_node/.style={font=\small}
}

\begin{tikzpicture}[x=0.65pt,y=0.65pt,yscale=-1,xscale=1]

% --- 1. 流形 M 与 N (严格保留原始路径) ---
\draw[manifold] (97,57.5) .. controls (117,47.5) and (205,20) .. (227,45.5) .. controls (249,71) and (254,138) .. (228,155.5) .. controls (202,173) and (111,184.5) .. (91,154.5) .. controls (71,124.5) and (77,67.5) .. (97,57.5) -- cycle ;
\draw[manifold] (376,59) .. controls (396,49) and (524,4.5) .. (546,39.5) .. controls (568,74.5) and (578,129.5) .. (548,166.5) .. controls (518,203.5) and (396,149) .. (376,119) .. controls (356,89) and (356,69) .. (376,59) -- cycle ;

% --- 2. 局部邻域 U 与 V (严格保留原始路径) ---
\draw[chart] (122,111) .. controls (122,91.95) and (140.13,76.5) .. (162.5,76.5) .. controls (184.87,76.5) and (203,111) .. (203,111) .. controls (203,130.05) and (184.87,145.5) .. (162.5,145.5) .. controls (140.13,145.5) and (122,130.05) .. (122,111) -- cycle ;
\draw[chart] (414,102.25) .. controls (414,83.33) and (436.83,68) .. (465,68) .. controls (493.17,68) and (516,83.33) .. (516,102.25) .. controls (516,121.17) and (493.17,136.5) .. (465,136.5) .. controls (436.83,136.5) and (414,121.17) .. (414,102.25) -- cycle ;

% --- 3. 欧氏空间中的像 (严格保留原始路径) ---
\draw[chart] (119.5,315.75) .. controls (119.5,289.38) and (140.88,268) .. (167.25,268) .. controls (193.62,268) and (215,289.38) .. (215,315.75) .. controls (215,342.12) and (193.62,363.5) .. (167.25,363.5) .. controls (140.88,363.5) and (119.5,342.12) .. (119.5,315.75) -- cycle ;
\draw[chart] (433.5,316.25) .. controls (433.5,290.71) and (454.21,270) .. (479.75,270) .. controls (505.29,270) and (526,290.71) .. (526,316.25) .. controls (526,341.79) and (505.29,362.5) .. (479.75,362.5) .. controls (454.21,362.5) and (433.5,341.79) .. (433.5,316.25) -- cycle ;

% --- 4. 点 p 和 f(p) (标准圆点) ---
\fill (162.5,111) circle (1.5pt) node[below right] {$p$};
\fill (465,102.25) circle (1.5pt) node[below right] {$f(p)$};

% --- 5. 调整后的水平/垂直箭头 ---

% 映射 f: 水平箭头 (调整 y 坐标使起点终点一致)
\draw[map_arrow] (240,105) -- (365,105) node[midway, above] {$f$};

% 坐标映射 phi: 垂直箭头 (调整 x 坐标使起点终点一致)
\draw[map_arrow] (167.25,150) -- (167.25,265) node[midway, left] {$\varphi$};

% 坐标映射 psi: 垂直箭头
\draw[map_arrow] (479.75,145) -- (479.75,267) node[midway, right] {$\psi$};

% 复合映射: 水平箭头
\draw[map_arrow] (220,316) -- (430,316) node[midway, below=5pt] {$\psi \circ f \circ \varphi^{-1}$};

% --- 6. 文本标注 ---
\node[label_node] at (105, 45) {$M$};
\node[label_node] at (565, 55) {$N$};
\node[label_node] at (145, 90) {$U$};
\node[label_node] at (445, 85) {$V$};
\node[label_node, anchor=north] at (167.25, 365) {$\varphi(U)$};
\node[label_node, anchor=north] at (479.75, 365) {$\psi(V)$};

\end{tikzpicture}
\end{center}
\begin{Definition}{微分同胚}
    设$M,N$为$C^{r}$微分流形,$f:M\to N$为同胚映射.如果$f$及其逆映射$f^{-1}$均为$C^{r}$映射,则称$f$为$C^{r}$微分同胚.
\end{Definition}
我们知道.流形其实就是局部看起来像欧氏空间的拓扑空间.设$\{(U_{\alpha},\varphi_{\alpha})\}$为微分流形$M$的局部坐标覆盖.转换映射可以用分量表示如下
$$
\varphi_{\beta}\circ\varphi^{-1}_{\alpha}(x)=(y^1,y^2,\cdots,y^n)
$$
其中$x=(x^1,x^2,\cdots,x^n)\in \varphi_{\alpha}(U_{\alpha}\cap U_{\beta}),y^i(1\leqslant i\leqslant n)$是$x^j(1\leqslant j\leqslant n)$的函数.于是转换映射$\varphi_{\beta}\circ\varphi^{-1}_{\alpha}$的Jacobi矩阵记为
$$
J(\varphi_{\beta}\circ\varphi^{-1}_{\alpha})=\left(\frac{\partial y^i}{\partial x^j}\right)_{n\times n}
$$
相应的Jaocbi行列式为
$$
\det J(\varphi_{\beta}\circ\varphi^{-1}_{\alpha})(p)=\frac{\partial (y^1,y^2,\cdots,y^n)}{\partial (x^1,x^2,\cdots,x^n)}(\varphi_{\alpha}(p)),\quad p\in U_{\alpha}\cap U_{\beta}
$$
为了后面在流形上定义积分,我们需要定向的概念.
\begin{Definition}{可定向流形}
    设$M$为
\end{Definition}
























考虑单位圆周
$$
S^1=\{(x,y)\in\mathbb{R}^2:x^2+y^2=1\}=\{\mathrm{e}^{\mathrm{i}\theta}:\theta\in[0,2\pi]\}
$$
则$S^1$为$\mathbb{R}^2$的子拓扑空间.令
$$
U_1=\{\mathrm{e}^{\mathrm{i}\theta}:\theta\in(0,2\pi)\},\quad U_2=\{\mathrm{e}^{\mathrm{i}\theta}:\theta\in(\pi,3\pi)\}
$$
则$U_1=S^1\setminus\{(1,0)\},U_2=S^1\setminus\{(-1,0)\}$,且$S^1=U_1\cup U_2$.分别定义映射如下
$$
\begin{aligned}
    \varphi_1:U_1&\to (0,2\pi)\subset \mathbb{R}^1\\
    \mathrm{e}^{\mathrm{i}\theta}&\mapsto \theta
\end{aligned}
$$
$$
\begin{aligned}
    \varphi_2:U_2&\to (\pi,3\pi)\subset \mathbb{R}^1\\
    \mathrm{e}^{\mathrm{i}\eta}&\mapsto \eta
\end{aligned}
$$
则$\varphi_1,\varphi_2$都是同胚映射,且它们的转换映射为
$$
\varphi_2\circ \varphi_1^{-1}:\varphi_1(U_1\cap U_2)=(0,\pi)\cup(\pi,2\pi)\to \varphi_2(U_1\cap U_2)=(\pi,2\pi)\cup (2\pi,3\pi),
$$
形如
$$
\varphi_2\circ \varphi_1^{-1}=
\begin{cases}
    \theta+2\pi, & \theta\in(0,\pi)\\[3pt]
    \theta, & \theta\in(\pi,2\pi)
\end{cases}
$$


































